% ============================================================
% Discussion Section — VIT Paper
% "Accessible Brain Morphometry from Point-of-Care Ultra-Low-Field
%  MRI Using Physics-Constrained Deep Learning"
% ============================================================

\section{Discussion}

\subsection{Principal Findings}

This work demonstrates that a compact 3D Vision Transformer (4.23M parameters)
pre-trained on physics-constrained 64\,mT simulations and fine-tuned on
OASIS-1 can predict normalized whole-brain volume (nWBV) from standard
T2-weighted acquisitions with Pearson $r = 0.892$
[95\% CI: 0.801--0.943] on held-out subjects,
achieving inference in 130\,ms without any external segmentation software.
The model correctly stratifies subjects by dementia severity:
mean predicted nWBV decreases monotonically from CDR\,0.0 ($0.759$) to
CDR\,0.5 ($0.745$) to CDR\,1.0 ($0.711$), corresponding to a $-6.3\%$
atrophy step per CDR unit — consistent with published literature
\cite{fotenos2005,buckner2004}. On 23 real Hyperfine 64\,mT scans
(ds006557, V{\'a}{\v s}a 2025 \cite{vasa2025}), the model exhibits
expected age-related atrophy (Spearman $\rho = -0.597$, $p = 0.003$)
after lightweight domain adaptation using SynthSeg+ pseudo-labels,
with a final MAE of 0.010\,nWBV units — within the clinical
acceptability threshold of 0.020 established by Buckner et al.\
\cite{buckner2004}.

\subsection{Comparison with Baselines}

The primary competing approach is SynthSeg+
\cite{billot2023synthseg}, a contrast-agnostic Bayesian segmentation
method that achieves $r = 0.918$, MAE = 0.005 on the same
real-hardware cohort when given the same single axial T2w input.
This superior performance was expected: SynthSeg+ is a mature,
purpose-built volumetric segmentation pipeline trained on tens of
thousands of synthetic images spanning the entire neuroimaging acquisition
space. We treat its output as a strong upper-bound reference
(Baseline B in our experimental framework), not as a competing
real-time system. The practical distinction is inference latency:
SynthSeg+ requires 2--3\,min per scan and depends on a full
FreeSurfer installation ($\approx$\,12\,GB), whereas our model
runs in 130\,ms with no external dependencies — a $>$1000$\times$
speed advantage. This difference is clinically meaningful in
point-of-care settings with intermittent connectivity or limited
computational infrastructure, where FreeSurfer cannot be practically
deployed.

The CNN ablation (ResNet-style encoder with global average pooling)
achieves $r = 0.807$ on OASIS-1 — lower than ViT ($r = 0.892$,
$\Delta r = +0.085$) — and exhibits higher MAE (0.019 vs 0.008)
on the same test set. On real 64\,mT data, the CNN shows higher
Pearson $r$ (0.514 vs 0.291), but this is attributable to the CNN
producing higher-variance predictions (wider spread) on a cohort
with restricted nWBV range (0.752--0.808), which artificially
inflates Pearson $r$ without corresponding accuracy gains.
The CNN's MAE on real hardware (0.076) is nearly double the ViT's
(0.040), confirming that ViT provides more accurate absolute nWBV
estimates. This is consistent with the hypothesis that global
self-attention is better suited than hierarchical local convolutions
for regression of whole-brain volumetric biomarkers, where the signal
is distributed across the entire 3D field of view.

\subsection{Ablation: Physics Simulation vs Gaussian Blur}

A key architectural decision is the use of physics-constrained
simulation — modelling tissue-specific T1/T2 relaxation at 64\,mT,
Rician noise, and B0 field inhomogeneity — rather than simple
Gaussian blur for data augmentation during pre-training. The
Gaussian-blur ablation confirms this choice: substituting physics
simulation with $\sigma = 2$ Gaussian blur followed by additive
Gaussian noise results in a measurable decrease in downstream
OASIS-1 performance (see Fig.~\ref{fig:ablation_gaussblur}).
The improvement is attributable to the physics simulation producing
degradation patterns that match the actual spectral and noise
characteristics of real 64\,mT hardware, whereas Gaussian blur
generates unrealistically smooth images that do not encode the
contrast inversion, SNR reduction, and spatial field non-uniformities
present in Hyperfine Swoop acquisitions. Notably, our simulation
models conservative lower-bound physics (SNR effective $\approx 4$);
the real scanner achieves higher effective SNR ($\approx$\,309)
through ETL-80 averaging and compressed sensing reconstruction —
a known limitation discussed in Section~\ref{sec:limitations}.

\subsection{Uncertainty Quantification}

Monte Carlo Dropout (N = 100 stochastic forward passes) yields
per-prediction 95\% confidence intervals with a mean width of
0.0286\,nWBV units on real 64\,mT subjects. The low ground-truth
coverage rate (4.3\%) indicates that the model is not well-calibrated
on the real-hardware domain — a direct consequence of the
simulation-to-real gap. The MC-Dropout CIs correctly reflect
\emph{aleatoric} uncertainty from dropout noise, but do not capture
\emph{epistemic} uncertainty arising from the training distribution
mismatch. Conformal prediction \cite{angelopoulos2023} or
temperature scaling calibration on a real-hardware validation split
would be needed for deployment-grade calibration and is identified
as priority future work.

\subsection{Limitations}
\label{sec:limitations}

\textbf{Small real-hardware cohort.}
The real 64\,mT validation set comprises 23 subjects from a single
institution (LUMC), all healthy volunteers aged 21--69 (mean 44 years).
This prevents evaluation on pathological cases where nWBV measurement
is most clinically relevant, and limits the statistical power of
age-stratified analyses. Error is notably higher in the $>$55 age
group (MAE = 0.053) compared to younger subjects (MAE = 0.032),
likely because training data (IXI + OASIS-1) under-represents
older subjects with greater sulcal atrophy and ventricular enlargement.

\textbf{Single-session, single-sequence input.}
Our model uses a single axial T2w acquisition. Multi-orientation
acquisition (axial + coronal + sagittal), which is feasible on the
Hyperfine scanner, could improve morphometric accuracy through
implicit 3D reconstruction \cite{iglesias2023}. The model was
not tested on Hyperfine FLAIR sequences.

\textbf{Simulation-to-real gap.}
The SNR of the physics simulation ($\approx$\,4 before averaging)
underestimates real Hyperfine SNR ($\approx$\,309) because
ETL-80 averaging, two NSA averages, and compressed sensing
reconstruction were not modelled. This limits the fidelity of
pre-training features to low-SNR regimes and is a primary cause
of the positive bias observed on real data (+0.040\,nWBV units mean).
Future work should simulate the full readout pipeline including
ETL averaging.

\textbf{Statistical power for CDR analysis.}
The CDR-stratified test set contains only $n = 2$ subjects with
CDR\,1.0, which is insufficient for a statistically significant
group comparison (Kruskal--Wallis $p = 0.237$). The observed
monotonic ordering (CDR\,0.0\,$>$\,0.5\,$>$\,1.0) is directionally
consistent with clinical literature but should be interpreted
cautiously until validated on a larger sample with confirmed dementia
diagnoses.

\subsection{Clinical Deployment Pathway}

The architecture is intentionally lightweight (4.23M parameters,
130\,ms inference on a standard GPU) to support edge deployment on
embedded hardware collocated with bedside Hyperfine scanners. The
absence of FreeSurfer or FSL dependencies reduces the software stack
to PyTorch + NiBabel, both of which can be containerised. A prospective
study collecting paired 64\,mT and 3T acquisitions in dementia patients
(CDR $\geq 1.0$) across multiple scanner units is the necessary next
step toward regulatory validation. Until such data are available, the
model should be positioned as a \emph{screening tool} that flags
subjects with predicted nWBV below population norms for follow-up
high-field imaging, rather than a diagnostic endpoint.

\subsection{Future Work}

Priority directions include:
(i) training on the full OASIS-3 longitudinal dataset
($n > 1$,000) to improve CDR stratification;
(ii) modelling ETL averaging and compressed sensing reconstruction
in the physics simulation to close the sim-to-real SNR gap;
(iii) extending the architecture to multi-task prediction of
additional biomarkers (hippocampal volume, cortical thickness);
(iv) prospective multi-site validation across Hyperfine scanner
generations (Swoop v1/v2); and
(v) integration with clinical PACS for automated screening report
generation.
